\en{In this chapter we prove that $m\cdot\wp(u;L)$ is an algebraic integer of $\mathbb Z\mathopen{}\left[\frac 1 4 g_2(L); \frac 1 4 g_3(L)\right]\mathclose{}$ for all positive integers $m$ and for all $u\in\mathbb C-L$ with $m\cdot u\in L$.}%
\de{In diesem Kapitel beweisen wir, dass $m\cdot\wp(u;L)$ für alle natürlichen Zahlen $m\neq 0$ und für alle $u\in\mathbb C-L$ mit $m\cdot u\in L$ ganzalgebraisch in $\mathbb Z\mathopen{}\left[\frac 1 4 g_2(L); \frac 1 4 g_3(L)\right]\mathclose{}$ ist.}

\en{This appendix elaborates \cite[pp.~184-185]{\en{EN}Fricke2}.}%
\de{Dieser Anhang arbeitet \cite[pp.~184-185]{\en{EN}Fricke2} aus.}

\begin{defi}\label{\en{EN}defidivi}
\en{Given a complex number $m\neq 0$ and a lattice $L$.
Then $u$ is called a $m$-division-point and $\wp(u)$ is called a $m$-division value iff it holds:}%
\de{Gegeben sind eine komplexe Zahl $m\neq 0$ und ein Gitter $L$.
Dann nennen wir $u$ eine $m$-Teilungsstelle und $\wp(u)$ einen $m$-Teilungswert genau dann, wenn:}
$$m\cdot u\in L,\quad\text{ \en{but}\de{aber} }\quad u\notin L.$$
\en{We denote the following set $\DIV(m)$, which contains all $m$-division-points in the fundamental parallelogram $\mathcal P$ from Def.~\ref{\en{EN}fund}:}%
\de{Wir bezeichnen die folgende Menge mit $\DIV(m)$, sie enthält alle $m$-Teilungsstellen im Fundamentalparallelogramm $\mathcal P$ aus Def.~\ref{\en{EN}fund}:}
\begin{align*}
\DIV(m)&=\left\{~u\in\mathcal P~\left|~m\cdot u\in L,\text{ \en{but }\de{aber } }u\notin L\right.\right\}\\
\text{\en{where}\de{wobei}}\quad \mathcal P&=\left\{~a\omega_1+b\omega_2\in\mathbb C~\left|~0\leq a,b < 1\right.\right\}
\end{align*}
\en{Throughout Appendix~\ref{\en{EN}kapdivi}, $m$ will even be a positive integer.

\noindent Fig.~\ref{\en{EN}abbwege} on p.~\pageref{\en{EN}abbwege} shows $\overline{\mathcal P}$ and those points in $\overline{\mathcal P}$ which are equivalent to $u\in \DIV(2)$.}%
\de{Im gesamten Anhang~\ref{\en{EN}kapdivi} wird $m$ sogar positiv ganzzahlig sein.

\noindent Abb.~\ref{\en{EN}abbwege} auf S.~\pageref{\en{EN}abbwege} zeigt $\overline{\mathcal P}$ und die Stellen in $\overline{\mathcal P}$, die zu $u\in \DIV(2)$ äquivalent sind.}%
\end{defi}


\begin{lem}\label{\en{EN}lemmm}
\en{If $m\neq 0$ is integral, then the number of $m$-division-points in $\mathcal P$ is $m^2-1$.}%
\de{Wenn $m\neq 0$ ganzzahlig ist, dann gibt es $m^2-1$ $m$-Teilungs-Punkte in $\mathcal P$.}
\end{lem}
\begin{proof}
\en{$u\in \DIV(m)$ is equivalent to $m\cdot u = k\omega_1+l\omega_2$ with integral $k,l$ such that $u=\frac{k}{m}\cdot \omega_1+\frac{l}{m}\cdot\omega_2\in \mathcal P-\{0\}$.
So every pair $(k,l)\in\mathbb Z^2-\{(0,0)\}$ with $0\leq k,l < m$ yields one $u$. Since there are $m^2-1$ such pairs, the Lemma is proven.}%
\de{$u\in \DIV(m)$ ist äquivalent zu $m\cdot u = k\omega_1+l\omega_2$ mit ganzzahligen $k,l$ so dass $u=\frac{k}{m}\cdot \omega_1+\frac{l}{m}\cdot\omega_2\in \mathcal P-\{0\}$.
Also liefert jedes Paar $(k,l)\in\mathbb Z^2-\{(0,0)\}$ mit $0\leq k,l < m$ eine Stelle $u$. Weil es $m^2-1$ solche Paare gibt, ist das Lemma bewiesen.}
\end{proof}


\begin{lem}\label{\en{EN}lemdefifm}
\en{The following quotient of Weierstraß $\sigma$-functions (cf.~Def.~\ref{\en{EN}defisigma}) is an elliptic function with periods $\omega_{1;2}$:}%
\de{Der folgende Quotient Weierstraß'scher $\sigma$-Funktionen (Def.~\ref{\en{EN}defisigma}) ist eine elliptische Funktion mit Perioden $\omega_{1;2}$:}
$$F_m(z):=\frac{\sigma(m\cdot z)}{(\sigma(z))^{m\cdot m}}$$
\en{It has a pole of order $m^2-1$ at $z\in L$ (e.g.~at $z=0$), and $m^2-1$ zeros of order one at $z\in \DIV(m)$.
Modulo $L$, $F_m$ has no further poles or zeros.}%
\de{Sie hat bei $z\in L$ (z.B.~bei $z=0$) einen Pol der Ordnung $m^2-1$ und $m^2-1$ einfache Nullstellen bei $z\in \DIV(m)$.
Modulo $L$ hat $F_m$ keine weiteren Null- oder Polstellen.}
\end{lem}
\begin{proof}
\en{In Prop.~\ref{\en{EN}trafosigma} we proved}%
\de{Satz~\ref{\en{EN}trafosigma} besagt}
$$\sigma(z+\omega_k)=-\exp\lk\eta_k\cdot\left(z+\frac{\omega_k}{2}\right)\rk\cdot \sigma(z)$$
\en{From this we will prove by induction that}%
\de{Hieraus werden wir per vollständiger Induktion beweisen:}
\begin{align}
    \sigma(z+n\omega_k) &= (-1)^n \cdot \exp\lk n\cdot\eta_k\left(z+n\cdot\frac{\omega_k}{2}\right)\rk \cdot ~\sigma(z)\label{\en{EN}glgsigmatr}
\end{align}
\en{For $n=0$ we have $\sigma(z)=\sigma(z)$. Then we do the step $n-1\rightarrow n$, where we first use Prop.~\ref{\en{EN}trafosigma}:}%
\de{Zum Induktionsanfang $n=0$: Hier besagt Glg.~(\ref{\en{EN}glgsigmatr}) nur, dass $\sigma(z)=\sigma(z)$ ist.\\
Im Induktionsschritt $n-1\rightarrow n$ nutzen wir zunächst Satz~\ref{\en{EN}trafosigma}:}
\begin{align*}
    \sigma(z+n\omega_k) &= \sigma((z+(n-1)\omega_k)+~\omega_k)\\
    &= \underbrace{- \exp\lk \eta_k\left(z+(n-1)\omega_k+\frac{\omega_k}{2}\right)\rk}_{P} \cdot ~\sigma(z+(n-1)\omega_k)
\end{align*}
\en{Then we use the induction basis:}%
\de{Dann benutzen wir die Induktionsvoraussetzung:}
\begin{align*}
\sigma(z+(n-1)\omega_k) = \underbrace{(-1)^{n-1} \cdot \exp\lk (n-1)\cdot\eta_k\left(z+(n-1)\cdot\frac{\omega_k}{2}\right)\rk}_Q \cdot~\sigma(z)    
\end{align*}
\en{These factors can be combined to:}%
\de{Dann fassen wir diese Faktoren zusammen:}
\begin{align*}
  P\cdot Q &= (-1)^n \cdot \exp\lk \eta_k \cdot\left( z+(n-1)\omega_k+\frac{\omega_k}{2} + (n-1)\cdot\left(z+(n-1)\cdot\frac{\omega_k}{2}\right)   \right)\rk\\
  &= (-1)^n \cdot \exp\lk \eta_k \cdot\left(n\cdot z+\omega_k\cdot\left(n-1 + \frac 1 2 + \frac{(n-1)^2}{2}\right)   \right)\rk\\
  &= (-1)^n \cdot \exp\lk \eta_k \cdot\left(n\cdot z+\frac{\omega_k}{2}\cdot n^2   \right)\rk
\end{align*}
\en{This proves~(\ref{\en{EN}glgsigmatr}).
Now we prove the periodicity of $F_m$ using~(\ref{\en{EN}glgsigmatr}):}%
\de{Somit ist~(\ref{\en{EN}glgsigmatr}) bewiesen.
Dies benutzen wir nun, um die Perioden von $F_m$ zu beweisen:}
\begin{align*}
    F_m(z+\omega_k) &= \frac{\sigma(m\cdot z+m\omega_k)}{(\sigma(z+\omega_k))^{m\cdot m}}
    = \frac{(-1)^m \cdot \exp\lk m\cdot\eta_k\left(m\cdot z+m\cdot\frac{\omega_k}{2}\right)\rk \cdot ~\sigma(m\cdot z)}{\left(-\exp\lk\eta_k\cdot\left(z+\frac{\omega_k}{2}\right)\rk\cdot \sigma(z)\right)^{m\cdot m}}
\end{align*}
\en{But for all integral $m$ we have $(-1)^m=(-1)^{m\cdot m}$ and}%
\de{Aber für alle ganzzahligen $m$ gilt $(-1)^m=(-1)^{m\cdot m}$ und}
$$\exp\lk m\cdot\eta_k\left(m\cdot z+m\cdot\frac{\omega_k}{2}\right)\rk = \left(\exp\lk \eta_k\left(z+\frac{\omega_k}{2}\right)\rk\right)^{m\cdot m},$$
\en{thus we have proven $F_m(z+\omega_k)=F_m(z)$.
Since the $\sigma$-function has zeros of order one in all points of the lattice (cf.~Def.~\ref{\en{EN}defisigma}), the nominator is zero iff $mz\in L$ and the denominator is zero iff $z\in L$.
This proves that $F_m(z)$ has zeros of order one for all $z\in \DIV(m)$ and that $F_m$ has a pole of order $m^2-1$ at $z\in L$.}%
\de{somit haben wir $F_m(z+\omega_k)=F_m(z)$ bewiesen.
Da die $\sigma$-Funktion in allen Gitterpunkten einfache Nullstellen hat (siehe Def.~\ref{\en{EN}defisigma}), liefern genau diejenigen $z$ mit $mz\in L$ Nullstellen des Zählers und die mit $z\in L$ Nullstellen des Nenners.
Somit hat $F_m(z)$ einfache Nullstellen bei $z\in \DIV(m)$ und einen Pol der Ordnung $m^2-1$ bei $z\in L$.}
\end{proof}

\begin{lem}\label{\en{EN}lemnullst}
\en{For any positive integer $m$, the function}%
\de{Für alle natürlichen $m\neq 0$ hat die Funktion}
$$h_m(z) := m^2\cdot\prod_{u\in \DIV(m)}\left(\wp(z)-\wp(u)\right)$$
\en{has zeros of order two for all $z\in \DIV(m)$ and no further zeros modulo $L$.}%
\de{doppelte Nullstellen für alle $z\in \DIV(m)$ und keine weiteren Nullstellen modulo $L$.}%
\end{lem}
\begin{proof}
\en{If $z\in \DIV(m)$ and $z\notin \DIV(2)$, we deduce $2z\not\equiv 0$ and $z\not\equiv -z$.
Prop.~\ref{\en{EN}pgerade} tells $\wp(-z)=\wp(z)$, thus both the factor with $u\equiv z$ and the factor with $u\equiv -z$ vanish and we get a zero of order two for these $z$.}%
\de{Falls $z\in \DIV(m)$ und $z\notin \DIV(2)$ gilt, folgt $2z\not\equiv 0$ und $z\not\equiv -z$.
Dann besagt Satz~\ref{\en{EN}pgerade}, dass $\wp(-z)=\wp(z)$ gilt und folglich sowohl der Faktor mit $u\equiv z$ als auch der mit $u\equiv -z$ verschwindet, weshalb wir bei solchen $z$ eine \emph{doppelte} Nullstelle haben.}

\en{If there are $z\in \DIV(m)\cap \DIV(2)$, we combine these three factors like in Prop.~\ref{\en{EN}pstrichprod}:
$\left(\wp(z)-e_1\right) \cdot\left(\wp(z)-e_2\right) \cdot\left(\wp(z)-e_3\right)=\frac 1 4 \wp'(z)^2$.
From the zeros of $\wp'$ in Prop.~\ref{\en{EN}zerowp} we see that $h_m$ has zeros of order two for these $z$ as well.}%
\de{Falls noch $z\in \DIV(m)\cap \DIV(2)$ verbleiben, fassen wir diese drei Faktoren wie in Satz~\ref{\en{EN}pstrichprod} zusammen:
$\left(\wp(z)-e_1\right) \cdot\left(\wp(z)-e_2\right) \cdot\left(\wp(z)-e_3\right)=\frac 1 4 \wp'(z)^2$.
Aus den Nullstellen von $\wp'$ in Satz~\ref{\en{EN}zerowp} folgt, dass $h_m$ auch bei diesen $z$ doppelte Nullstellen hat.}


\en{Since there are $m^2-1$ factors in $h_m(z)$ (Lemma~\ref{\en{EN}lemmm}) and $\wp(z)$ has the order two (Def.~\ref{\en{EN}defiwp}), $h_m$ has the order $2\cdot(m^2-1)$.
Thus (by the third Liouville theorem, Prop.~\ref{\en{EN}liouville3}) $h_m$ has no further zeros modulo $L$.}%
\de{Da $m^2-1$ Faktoren in $h_m(z)$ sind (Lemma~\ref{\en{EN}lemmm}) und $\wp(z)$ die Ordnung zwei hat (Def.~\ref{\en{EN}defiwp}), hat $h_m$ die Ordnung $2\cdot(m^2-1)$.
Aus dem dritten Liouville'schen Satz~\ref{\en{EN}liouville3} folgt daher, dass $h_m$ modulo $L$ keine weiteren Nullstellen hat.}
\end{proof}

\begin{thm}\label{\en{EN}propfmprod}
\en{For the $F_m$ and $h_m$ from Lemma~\ref{\en{EN}lemdefifm} and \ref{\en{EN}lemnullst}, it holds $F_m^2(z)=h_m(z)$.}%
\de{Für die $F_m$ aus Lemma~\ref{\en{EN}lemdefifm} und die $h_m$ aus Lemma~\ref{\en{EN}lemnullst} gilt $F_m^2(z)=h_m(z)$.}
\end{thm}
\begin{proof}
\en{Both $F_m^2(z)$ and $h_m(z)$ are elliptic functions with no poles outside of $L$.
From $\sigma(z)\approx z$ (Def.~\ref{\en{EN}defisigma}) and from Lemma~\ref{\en{EN}lemdefifm} we see that the Laurent series of $F_m(z)^2$ starts with $m^2\cdot z^{-2(m^2-1)}$. The same is true for $h_m(z)$ (cf.~Lemma~\ref{\en{EN}lemmm} and $\wp(z)\approx 1/z^2$ from Prop.~\ref{\en{EN}laurentwp}), thus the quotient $q_m:=h_m/F_m^2$ takes the value $1$ at $z=0$, which proves that the quotient $q_m:=h_m/F_m^2$ has no pole at $z\in L$.}%
\de{Sowohl $F_m^2(z)$ als auch $h_m(z)$ sind elliptische Funktionen, die keine Pole außerhalb von $L$ haben.
Aus $\sigma(z)\approx z$ (Def.~\ref{\en{EN}defisigma}) und aus Lemma~\ref{\en{EN}lemdefifm} folgt, dass die Laurentreihe von $F_m(z)^2$ mit $m^2\cdot z^{-2(m^2-1)}$ beginnt. $h_m(z)$ beginnt ebenso (wegen $\wp(z)\approx 1/z^2$ aus Satz~\ref{\en{EN}laurentwp} und wegen Lemma~\ref{\en{EN}lemmm}), folglich hat der Quotient $q_m:=h_m/F_m^2$ bei $z=0$ den Wert $1$ (und keinen Pol bei $z\in L$).}

\en{The quotient $q_m$ could still have poles in the zeros of $F_m^2$, but since $F_m^2$ and $h_m$ have the same zeros (Lemma~\ref{\en{EN}lemdefifm} and~\ref{\en{EN}lemnullst}), $q_m$ is elliptic without poles. By the first Liouville theorem (Prop.~\ref{\en{EN}liouville1}), the quotient is constant.

We just proved $q_m(0)=1$, thus we obtain $q_m(z)=1$ and $F_m^2(z)=h_m(z)$.}%
\de{Der Quotient $q_m$ könnte noch Pole in den Nullstellen von $F_m^2$ haben, aber da $F_m^2$ und $h_m$ die gleichen Nullstellen haben (Lemma~\ref{\en{EN}lemdefifm} und~\ref{\en{EN}lemnullst}), ist $q_m$ eine elliptische Funktion ohne Pole. Aufgrund des ersten Liouville'schen Satzes~\ref{\en{EN}liouville1} ist der Quotient konstant.

Soeben haben wir $q_m(0)=1$ bewiesen, also gilt $q_m(z)=1$ und $F_m^2(z)=h_m(z)$.}
\end{proof}

\begin{bem}
\en{Prop.~\ref{\en{EN}propfmprod} gives $F_m^2(z)$ as a polynomial in $\wp(z)$. In the rest of this chapter, we will construct this polynomial recursively and use this recursion to prove that $m\cdot\wp(u)$ is an algebraic integer of~\,$\mathbb Z\mathopen{}\left[\frac 1 4 g_2(L);\frac 1 4 g_3(L)\right]\mathclose{}$.}%
\de{Satz~\ref{\en{EN}propfmprod} stellt $F_m^2(z)$ als Polynom in $\wp(z)$ dar. Im Rest dieses Kapitels werden wir dieses Polynom rekursiv konstruieren und diese Rekursion nutzen, um zu beweisen, dass $m\cdot\wp(u)$ für $u\in \DIV(m)$ ganzalgebraisch in $\mathbb Z\mathopen{}\left[\frac 1 4 g_2(L);\frac 1 4 g_3(L)\right]\mathclose{}$ ist.}
\end{bem}


\begin{lem}\label{\en{EN}lemwpsigma}
\en{For all $u,v\notin L$ it holds}%
\de{Für alle $u,v\notin L$ gilt}
$$\wp(v) - \wp (u) = \frac{\sigma(u+v)\sigma(u-v)}{\sigma^2(u)\sigma^2(v)}$$
\end{lem}
\begin{proof}
\en{We define the function}%
\de{Wir definieren die Funktion}
$$F(u,v):=\frac{\sigma(u+v)\sigma(u-v)}{\sigma^2(u)\sigma^2(v)}+\wp(u)-\wp(v)$$
\en{First we fix $v\notin L$ and analyze $F(u,v)$ as a function $g(u)$ which has no poles outside $L$. Around $u=0$, we use $\sigma(u)\approx u$ (cf.~Def.~\ref{\en{EN}defisigma}) and obtain:}%
\de{Zunächst fixieren wir $v\notin L$ und betrachten $F(u,v)$ als Funktion $g(u)$, welche keine Pole außerhalb $L$ hat. Um $u=0$ nutzen wir $\sigma(u)\approx u$ (vgl.~Def.~\ref{\en{EN}defisigma}) und erhalten:}
$$g(u) \approx\frac{\sigma(v)\sigma(-v)}{u^2\sigma^2(v)}+\frac{1}{u^2}-\wp(v)$$
\en{But since $\sigma(-v)=-\sigma(v)$ (cf.~Prop.~\ref{\en{EN}sigmaodd}), the two terms $\pm u^{-2}$ in the Laurent series of $g(u)$ cancel each other out.
And since $g(-u)=g(u)$, there can't be a pole of order one at $u=0$, which proves that $g(u)$ has no poles.

Prop.~\ref{\en{EN}trafosigma} shows that $g(u)$ is elliptic (details on this can be found in the calculation of eq.~(\ref{\en{EN}eqellip}) on p.~\pageref{\en{EN}eqellip}), thus $g(u)=F(u,v)$ is constant with respect to $u$.
In the same way we can prove that $F(u,v)$ is constant with respect to $v$, thus $F(u,v)$ is constant.
From Def.~\ref{\en{EN}defisigma} we get $\sigma(0)=0$ and thus $F(v,v)=0$, so $F(u,v)=0$ for all $u,v\notin L$.}%
\de{Satz~\ref{\en{EN}sigmaodd} besagt, dass $\sigma(-v)=-\sigma(v)$ gilt, also heben sich die beiden Terme $\pm u^{-2}$ in der Laurentreihe von $g(u)$ gegenseitig auf.
Weil weiter $g(-u)=g(u)$ gilt, kann es keinen Pol der Ordnung $1$ bei $u=0$ geben. Somit ist bewiesen, dass $g(u)$ keine Pole hat.

Aus Satz~\ref{\en{EN}trafosigma} folgt, dass $g(u)$ elliptisch ist (siehe hierzu auch die Rechnung zu Glg.~(\ref{\en{EN}eqellip}) auf S.~\pageref{\en{EN}eqellip}), also ist $g(u)=F(u,v)$ konstant bezüglich $u$.
Gleichermaßen können wir beweisen, dass $F(u,v)$ konstant bezüglich $v$ ist, also ist $F(u,v)$ konstant.
Aus Def.~\ref{\en{EN}defisigma} folgt $\sigma(0)=0$ und somit $F(v,v)=0$, also ist $F(u,v)=0$ für alle $u,v\notin L$.}
\end{proof}

\begin{lem}\label{\en{EN}lemwpnz}
\en{Using the $F_m$ from Lemma~\ref{\en{EN}lemdefifm}, it holds}%
\de{Für die $F_m$ aus Lemma~\ref{\en{EN}lemdefifm} gilt}
$$\wp(nz)= \wp(z) - \frac{F_{n-1}(z)\cdot F_{n+1}(z)}{F_n(z)^2}$$
\end{lem}
\begin{proof}
\en{Using Lemma~\ref{\en{EN}lemwpsigma} with $u=z$ and $v=nz$ yields:}%
\de{Lemma~\ref{\en{EN}lemwpsigma} liefert mit $u=z$ und $v=nz$:}
$$\wp(nz) - \wp (z) = \frac{\sigma((n+1)z)\sigma(-(n-1)z)}{\sigma^2(z)\sigma^2(nz)}$$
\en{Next we use $\sigma(z)=-\sigma(-z)$ (Prop.~\ref{\en{EN}sigmaodd}) and the definition of $F_k$ from Lemma~\ref{\en{EN}lemdefifm} which yields $\sigma(kz)=F_k(z)\cdot \sigma(z)^{k\cdot k}$:}
\de{Dann nutzen wir $\sigma(z)=-\sigma(-z)$ (Satz~\ref{\en{EN}sigmaodd}) und die Definition der $F_k$ aus Lemma~\ref{\en{EN}lemdefifm} in der Form $\sigma(kz)=F_k(z)\cdot \sigma(z)^{k\cdot k}$:}
\begin{align*}
    \wp(nz) - \wp (z) &= - \frac{F_{n+1}(z)\sigma(z)^{(n+1)\cdot(n+1)}\cdot F_{n-1}(z)\sigma(z)^{(n-1)\cdot(n-1)}}{\sigma^2(z)\cdot F_n(z)^2\sigma(z)^{2\cdot n\cdot n}}
\end{align*}
\en{Reducing this fraction by $\sigma(z)^{2n^2+2}$ proves the Lemma.}%
\de{Indem wir diesen Bruch mit $\sigma(z)^{2n^2+2}$ kürzen, ist das Lemma bewiesen.}
\end{proof}

\begin{lem}\label{\en{EN}lemwp2str}
\en{For all $z\in\mathbb C$ it holds}%
\de{Für alle $z\in\mathbb C$ gilt}
$\wp''(z) = 6\wp(z)^2-\frac 1 2 g_2$
\en{and}\de{und}
$$\wp(2z) = \frac 1 4 \cdot \left(\frac{\wp''(z)}{\wp'(z)}\right)^2-2\wp(z)$$
\end{lem}
\begin{proof}
\en{The algebraic differential equation of the $\wp$-function (Prop.~\ref{\en{EN}dglP}) yields:}%
\de{Die algebraische Differentialgleichung der $\wp$-Funktion aus Satz~\ref{\en{EN}dglP} lautet:}
\begin{align*}
    \wp'(z)^2 &= 4\wp(z)^3-g_2\wp(z)-g_3\qquad\left|\frac{d}{dz}\right.\\
    \Longrightarrow\quad 2\wp'(z)\wp''(z) &= 12 \wp(z)^2\wp'(z) - g_2\wp'(z)
\end{align*}
\en{Dividing by $2\cdot\wp'(z)$ yields $\wp''(z)$.}%
\de{Indem wir dies durch $2\cdot\wp'(z)$ dividieren, erhalten wir $\wp''(z)$.}

\en{Next we recall the Laurent series of $\wp(z)$ from Prop.~\ref{\en{EN}laurentwp}:}%
\de{Als nächstes rufen wir uns die Laurentreihe von $\wp(z)$ aus Satz~\ref{\en{EN}laurentwp} ins Gedächtnis:}
\begin{align*}
    \wp(z) &= z^{-2} + 3 G_4 z^2  + 5 G_6 z^4 + 7 G_8 z^6  + O(z^8)\\
    \wp'(z) &= -2z^{-3} + 6 G_4 z + 20 G_6 z^3 + 42 G_8 z^5  + O(z^7)\\
    \wp''(z) &= 6z^{-4} + 6 G_4   + 60 G_6 z^2 + 210 G_8 z^4  + O(z^6)
\end{align*}
\en{If we put this into $\wp''(z) = 6\wp(z)^2-\frac 1 2 g_2$ and compare the coefficients of $z^4$, we obtain:}%
\de{Wenn wir diese Reihen in $\wp''(z) = 6\wp(z)^2-\frac 1 2 g_2$ einsetzen und die Koeffizienten vor $z^4$ vergleichen, erhalten wir:}
$$210 G_8 = 6\cdot\left(2\cdot 7 G_8 + 9 G_4^2\right)\qquad\Longrightarrow\qquad 7G_8 = 3 G_4^2$$
\en{To prove the second equation, compare the Laurent series expansions of}%
\de{Um die zweite Gleichung zu beweisen, vergleichen wir die Laurentreihen von}
$\left(\frac{\wp''(z)}{2}\right)^2$ \en{and}\de{und} $\wp'(z)^2\cdot\left(\wp(2z)+2\wp(z)\right)$ \en{using}\de{unter Nutzung von} $7G_8=3 G_4^2$:
\begin{align*}
    \wp(z) &= z^{-2} + 3 G_4 z^2  + 5 G_6 z^4 + 3 G_4^2 z^6  + O(z^8)\\
    \wp'(z) &= -2z^{-3} + 6 G_4 z + 20 G_6 z^3 + 18 G_4^2 z^5  + O(z^7)\\
    \wp''(z) &= 6z^{-4} + 6 G_4   + 60 G_6 z^2 + 90 G_4^2 z^4  + O(z^6)\\
    \wp''(z)/2 &= 3z^{-4} + 3 G_4   + 30 G_6 z^2 + 45 G_4^2 z^4  + O(z^6)\\
    (\wp''(z)/2)^2 &= 9z^{-8} + 18 G_4 z^{-4} + 180 G_6 z^{-2} + 279 G_4^2 + O(z^2)
    \end{align*}\begin{align*}
    \wp'(z)^2 &= 4 z^{-6} - 24 G_4 z^{-2} - 80 G_6 - 36 G_4^2 z^2 + O(z^4)\\
    \wp(2z) &= 0\ko25 z^{-2} + 12 G_4 z^2  + 80 G_6 z^4 + 192 G_4^2 z^6  + O(z^8)\\
    \wp(2z)+2\wp(z) &= 2\ko25 z^{-2} + 18 G_4 z^2  + 90 G_6 z^4 + 198 G_4^2 z^6  + O(z^8)\\
    \wp'(z)^2\cdot\left(\wp(2z)+2\wp(z)\right) &= 9z^{-8} + 18 G_4 z^{-4} + 180 G_6 z^{-2} + 279 G_4^2 + O(z^2)
\end{align*}
\en{Since both $\left(\frac{\wp''(z)}{2}\right)^2$ and $\wp'(z)^2\cdot\left(\wp(2z)+2\wp(z)\right)$ are elliptic function without poles outside of $L$ (because the additional poles of $\wp(2z)$ are eliminated by the zeros of $\wp'(z)^2$ from Prop.~\ref{\en{EN}zerowp}) and their Laurent series are equal up to $O(z^2)$, the first Liouville theorem~\ref{\en{EN}liouville1} yields:}%
\de{Sowohl $\left(\frac{\wp''(z)}{2}\right)^2$ als auch $\wp'(z)^2\cdot\left(\wp(2z)+2\wp(z)\right)$ sind elliptische Funktionen, die keine Pole außerhalb von $L$ haben (weil die zusätzlichen Pole von $\wp(2z)$ durch die Nullstellen von $\wp'(z)^2$ aus Satz~\ref{\en{EN}zerowp} aufgehoben werden). Außerdem stimmen ihre Laurentreihen bis auf $O(z^2)$ überein, also liefert der erste Liouville'sche Satz~\ref{\en{EN}liouville1}:}
\begin{align*}
    \left(\frac{\wp''(z)}{2}\right)^2 - \wp'(z)^2\cdot\left(\wp(2z)+2\wp(z)\right) = 0
\end{align*}
\en{Solving this equation for $\wp(2z)$ produces the equation of Lemma~\ref{\en{EN}lemwp2str}.}%
\de{Wenn wir diese Gleichung nach $\wp(2z)$ auflösen, erhalten wir die zu beweisende Gleichung des Lemmas~\ref{\en{EN}lemwp2str}.}
\end{proof}


\begin{lem}\label{\en{EN}lemsigmasigma}
\en{For all $u, u_1, u_2, u_3 \in \mathbb C$ it holds}%
\de{Für alle $u, u_1, u_2, u_3 \in \mathbb C$ gilt}
\begin{align*}
  \sigma(u+u_1)\sigma(u-u_1)\sigma(u_2+u_3)\sigma(u_2-u_3)&\\
+~\sigma(u+u_2)\sigma(u-u_2)\sigma(u_3+u_1)\sigma(u_3-u_1)&\\
+~\sigma(u+u_3)\sigma(u-u_3)\sigma(u_1+u_2)\sigma(u_1-u_2)& = 0
\end{align*}
\en{where $\sigma$ denotes the Weierstraß $\sigma$-function.}%
\de{wobei $\sigma$ die Weierstraß'sche $\sigma$-Funktion bezeichnet.}
\end{lem}
\begin{proof}
\en{We distinguish two cases:}\de{Wir unterscheiden zwei Fälle:}
\en{In the \emph{first case}, at least two of the four complex numbers $u, u_1, u_2, u_3$ are equal modulo $L$. But then the equation of the Lemma is true, because one of the three terms is zero (since $\sigma(0)=0$) and the other two cancel each other out (since $\sigma(-z)=-\sigma(z)$).
For example $u_2=u_3$ yields:}%
\de{Im \emph{ersten Fall} sind mindestens zwei der vier komplexen Zahlen $u, u_1, u_2, u_3$ äquivalent modulo $L$. Dann liefert $\sigma(0)=0$, dass einer der drei Summanden Null ist, und die anderen beiden heben sich gegenseitig auf (weil $\sigma(-z)=-\sigma(z)$ ist).
Beispielsweise liefert $u_2=u_3$:}
\begin{align*}
  \sigma(u+u_1)\sigma(u-u_1)\sigma(u_2+u_2)\sigma(u_2-u_2)&\\
+~\sigma(u+u_2)\sigma(u-u_2)\sigma(u_2+u_1)\sigma(u_2-u_1)&\\
+~\sigma(u+u_2)\sigma(u-u_2)\sigma(u_1+u_2)\sigma(u_1-u_2)&\\
= 0 + \sigma(u+u_2)\sigma(u-u_2)\sigma(u_2+u_1)\sigma(u_2-u_1)&\\
    -~\sigma(u+u_2)\sigma(u-u_2)\sigma(u_2+u_1)\sigma(u_2-u_1)&=0
\end{align*}

\en{In the \emph{second case}, we have four complex numbers $u$, $u_1$, $u_2$ and $u_3$ which are pairwise distinct modulo $L$. Then we choose $u_4$ so that $u_4$ is not equivalent to any of the numbers $\{\pm u; \pm u_1; \pm u_2; \pm u_3; -u_4\}$ modulo $L$.}%
\de{Im \emph{zweiten Fall} haben wir vier komplexe Zahlen $u$, $u_1$, $u_2$ und $u_3$ welche paarweise inäquivalent modulo $L$ sind. Dann wählen wir $u_4$ so, dass $u_4$ inäquivalent zu allen Zahlen $\{\pm u; \pm u_1; \pm u_2; \pm u_3; -u_4\}$ modulo $L$ ist.}
\en{Next we define for $l\in\{1;2;3\}$ the functions}%
\de{Dann definieren für $l\in\{1;2;3\}$ die Funktionen}
$$f_l(z):=\frac{\sigma(z+u_l)\sigma(z-u_l)}{\sigma(z+u_4)\sigma(z-u_4)}$$
\en{From Prop.~\ref{\en{EN}trafosigma} we deduce that the $f_l(z)$ are elliptic functions:}%
\de{Aus Satz~\ref{\en{EN}trafosigma} folgt, dass die $f_l(z)$ elliptisch sind:}
\begin{align}
    f_l(z+\omega_k) &= \frac{\sigma(z+u_l+\omega_k)\cdot\sigma(z-u_l+\omega_k)}{\sigma(z+u_4+\omega_k)\cdot\sigma(z-u_4+\omega_k)}\nonumber\\
    &= \frac{\exp\lk\eta_k\left(z+u_l+\omega_k/2\right)\rk\cdot\exp\lk\eta_k\left(z-u_l+\omega_k/2\right)\rk}
    {\exp\lk\eta_k\left(z+u_4+\omega_k/2\right)\rk\cdot\exp\lk\eta_k\left(z-u_4+\omega_k/2\right)\rk}\cdot f_l(z) = f_l(z)\label{\en{EN}eqellip}
\end{align}
\en{From our choice of $u_4$, we see that it has two poles of order one at $\pm u_4$ (we have chosen $u_4$ such that the zeros of the denominator are inequivalent modulo $L$), thus it is an elliptic function of order two.}%
\de{Aus unserer Wahl von $u_4$ folgt, dass die $f_l$ zwei Pole erster Ordnung bei $\pm u_4$ haben (wir haben $u_4$ so gewählt, dass die Nullstellen des Nenners inäquivalent sind), also ist $f_l$ eine elliptische Funktion der Ordnung zwei.}
\en{Next we define the function}%
\de{Nun definieren wir die Funktion}
\begin{align*}
    f(z):= f_1(z)\cdot\sigma(u_2+u_3)\sigma(u_2-u_3)&\\
+f_2(z)\cdot\sigma(u_3+u_1)\sigma(u_3-u_1)&\\
+f_3(z)\cdot\sigma(u_1+u_2)\sigma(u_1-u_2)&
\end{align*}
\en{Since this is a linear combination of elliptic functions, it is also an elliptic function.}%
\de{Diese ist eine elliptische Funktion, weil sie eine Linearkombination elliptischer Funktionen ist.}
\en{But if we use $z\in\{u_1;u_2;u_3\}$, we obtain $f(z)=0$ (as in the first case). Since we know that these three numbers are pairwise inequivalent modulo $L$, the third Liouville theorem (Prop.~\ref{\en{EN}liouville3}) tells us that $f(z)$ must be a constant function (since a non-constant elliptic function of order two would have only two zeros modulo $L$), thus $f(z)$ is equal to zero.}%
\de{Aber für $z\in\{u_1;u_2;u_3\}$ erhalten wir $f(z)=0$ (wie im ersten Fall). Diese drei Zahlen sind aber paarweise inäquivalent modulo $L$. Aus dem dritten Liouville'schen Satz~\ref{\en{EN}liouville3} folgt nun, dass $f(z)$ konstant sein muss (weil eine nichtkonstante elliptische Funktion vom Grad zwei nur zwei Nullstellen modulo $L$ hätte) und somit  ist $f(z)=0$.}

\en{Finally we use $z=u$ with the originally given $u$ we get $f(u)=0$. But because $u_4$ is not equivalent to $\pm u$, we can multiply this with $\sigma(u+u_4)\sigma(u-u_4)$ and obtain the equation from the Lemma.}%
\de{Schließlich setzen wir $z=u$ mit dem ursprünglich gegebenen $u$ ein und erhalten $f(u)=0$.
Weiter gilt $u_4\not\equiv \pm u$ und wir dürfen $f(u)=0$ mit $\sigma(u+u_4)\sigma(u-u_4)$ multiplizieren. Somit ist das Lemma bewiesen.}
\end{proof}


\begin{lem}\label{\en{EN}lemfrekur}
\en{For the functions $F_m(z)$ from Lemma~\ref{\en{EN}lemdefifm}, it holds:}%
\de{Für die Funktionen $F_m(z)$ aus Lemma~\ref{\en{EN}lemdefifm} gilt:}
\begin{align*}
    F_{2n+1}(z) &=  F_{n+2}(z)\cdot F_n(z)^3 - F_{n-1}(z)\cdot F_{n+1}(z)^3 \\
    F_{2n}(z)\cdot F_2(z) &=  F_n(z)\cdot \left( F_{n+2}(z)\cdot F_{n-1}(z)^2 - F_{n-2}(z)\cdot F_{n+1}(z)^2\right)
\end{align*}
\end{lem}
\begin{proof}
\en{Using Lemma~\ref{\en{EN}lemsigmasigma} with $u=0$, $u_1=z$, $u_2=n\cdot z$ and $u_3=-(n+1)\cdot z$
yields:}%
\de{Wir setzen $u=0$, $u_1=z$, $u_2=n\cdot z$ und $u_3=-(n+1)\cdot z$ in Lemma~\ref{\en{EN}lemsigmasigma} ein:}
\begin{align*}
  \sigma(z)\sigma(-z)\sigma(-z)\sigma((2n+1)z)&\\
+~\sigma(nz)\sigma(-nz)\sigma(-nz)\sigma(-(n+2)z)&\\
+~\sigma(-(n+1)z)\sigma((n+1)z)\sigma((n+1)z)\sigma(-(n-1)z)& = 0
\end{align*}
\en{This can be simplified using $\sigma(-z)=-\sigma(z)$ from Prop.~\ref{\en{EN}sigmaodd}:}%
\de{Dann nutzen wir $\sigma(-z)=-\sigma(z)$ aus Satz~\ref{\en{EN}sigmaodd}:}
\begin{align*}
  \sigma(z)^3\sigma((2n+1)z)&=
  \sigma(nz)^3\sigma((n+2)z)
 -\sigma((n+1)z)^3\sigma((n-1)z)
\end{align*}
\en{Next we use the definition of the $F_k$ (cf.~Lemma~\ref{\en{EN}lemdefifm}) which yields $\sigma(kz)=F_k(z)\cdot \sigma(z)^{k\cdot k}$:}%
\de{Dann folgt aus der Definition der $F_k$ (in Lemma~\ref{\en{EN}lemdefifm}), dass $\sigma(kz)=F_k(z)\cdot \sigma(z)^{k\cdot k}$ gilt:}
\begin{align*}
  &~\sigma(z)^3 F_{2n+1}(z)\sigma(z)^{(2n+1)\cdot(2n+1)}\\
  =&~F_n(z)^3\sigma(z)^{3\cdot n\cdot n}\cdot F_{n+2}(z)\sigma(z)^{(n+2)\cdot(n+2)}\\
  &-F_{n+1}(z)^3\sigma(z)^{3\cdot (n+1)\cdot (n+1)}\cdot F_{n-1}(z)\sigma(z)^{(n-1)\cdot (n-1)}
\end{align*}
\en{Dividing by $\sigma(z)^{4n^2+4n+4}$ yields the first equation.}%
\de{Die erste Gleichung folgt hieraus durch Kürzen mit $\sigma(z)^{4n^2+4n+4}$.}

\en{To prove the second equation, we use Lemma~\ref{\en{EN}lemsigmasigma} with $u=\frac 1 2 z$, $u_1=\frac 3 2 z$, $u_2=\frac 1 2 \cdot(2n-1)\cdot z$ and $u_3=- \frac 1 2 \cdot(2n+1)\cdot z$:}%
\de{Um die zweite Gleichung zu beweisen, setzen wir $u=\frac 1 2 z$, $u_1=\frac 3 2 z$, $u_2=\frac 1 2 \cdot(2n-1)\cdot z$ und $u_3=- \frac 1 2 \cdot(2n+1)\cdot z$ in Lemma~\ref{\en{EN}lemsigmasigma} ein:}
\begin{align*}
  \sigma(2z)\sigma(-z)\sigma(-z)\sigma(2nz)&\\
+~\sigma(nz)\sigma(-(n-1)z)\sigma(-(n-1)z)\sigma(-(n+2)z)&\\
+~\sigma(-nz)\sigma((n+1)z)\sigma((n+1)z)\sigma(-(n-2)z)& = 0
\end{align*}
\en{Again, this can be simplified using $\sigma(-z)=-\sigma(z)$ to:}%
\de{Dann nutzen wir wieder $\sigma(-z)=-\sigma(z)$:}
\begin{align*}
 &~\sigma(2z)\sigma(z)^2\sigma(2nz)\\
=&~\sigma(nz)\sigma((n-1)z)^2\sigma((n+2)z)\\
 &-\sigma(nz)\sigma((n+1)z)^2\sigma((n-2)z)
\end{align*}
\en{As above, we use the definition of the $F_k$ which yields $\sigma(kz)=F_k(z)\cdot \sigma(z)^{k\cdot k}$:}%
\de{Dann folgt wie oben aus der Definition der $F_k$, dass $\sigma(kz)=F_k(z)\cdot \sigma(z)^{k\cdot k}$ gilt und somit}
\begin{align*}
 &~F_2(z)\sigma(z)^{2\cdot 2}\cdot \sigma(z)^2\cdot F_{2n}(z)\sigma(z)^{2n\cdot 2n}\\
=&~F_n(z)\sigma(z)^{n\cdot n}\cdot F_{n-1}(z)^2\sigma(z)^{2\cdot(n-1)\cdot(n-1)} \cdot F_{n+2}(z)\sigma(z)^{(n+2)\cdot(n+2)}\\
 &-F_n(z)\sigma(z)^{n\cdot n} \cdot F_{n+1}(z)^2\sigma(z)^{2\cdot(n+1)\cdot(n+1)} \cdot F_{n-2}(z)\sigma(z)^{(n-2)\cdot(n-2)}
\end{align*}
\en{Dividing by $\sigma(z)^{4n^2+6}$ yields the second equation.}%
\de{Die zweite Gleichung folgt hieraus durch Kürzen mit $\sigma(z)^{4n^2+6}$.}
\end{proof}

\pagebreak

\begin{defi}\label{\en{EN}defipm}
\en{For any positive integer $m$, we define the polynomial $P_m(x)$ as follows:}%
\de{Für alle natürlichen $m\neq 0$ definieren wir das Polynom $P_m(x)$ wie folgt:}
\begin{align*}
    P_1 &= 1;\quad P_2 = 1;\quad P_3 = 3x^4-6h_2x^2-12h_3x-h_2^2\\
    P_4 &= 2x^6-10h_2x^4-40h_3x^3-10h_2^2x^2-8h_2h_3x-16h_3^2+2h_2^3\\
   \intertext{\en{and for}\de{und für} $k\geq 1$:}
   P_{4k+1} &= 16(x^3-h_2x-h_3)^2\cdot P_{2k+2}\cdot P_{2k}^3 - P_{2k-1}\cdot P_{2k+1}^3\\
    P_{4k+2} &= P_{2k+1}\cdot \left( P_{2k+3}\cdot P_{2k}^2 - P_{2k-1}\cdot P_{2k+2}^2\right)\\
    P_{4k+3} &= P_{2k+3}\cdot P_{2k+1}^3 - 16(x^3-h_2x-h_3)^2\cdot P_{2k}\cdot P_{2k+2}^3\\
    P_{4k+4} &= P_{2k+2}\cdot \left( P_{2k+4}\cdot P_{2k+1}^2 - P_{2k}\cdot P_{2k+3}^2\right)
\end{align*}
\end{defi}

\begin{thm}\label{\en{EN}propfmpm}
\en{For any positive integer $m$, it holds}%
\de{Für alle natürlichen $m\neq 0$ gilt}
\begin{align*}
    F_m(z) &= \cas{-\wp'(z) \cdot P_m(\wp(z))}{P_m(\wp(z))}
\end{align*}
\en{where $\wp(z)$ denotes the Weierstraß $\wp$-function of a lattice $L$ with $h_2=\frac 1 4 g_2(L)$ and $h_3 = \frac 1 4 g_3(L)$.}%
\de{Hierbei bezeichne $\wp(z)$ die Weierstraß'sche $\wp$-Funktion eines Gitters $L$ mit $h_2=\frac 1 4 g_2(L)$ und $h_3 = \frac 1 4 g_3(L)$.}
\end{thm}
\begin{proof}
\en{Throughout the proof, we will use the abbreviations $x=\wp(z)$, $h_2 = \frac 1 4 g_2(L)$ and $h_3 = \frac 1 4 g_3(L)$.
First, we prove $m\leq 4$ as induction basis:}%
\de{Während des gesamten Beweises nutzen wir die Abkürzungen $x=\wp(z)$, $h_2 = \frac 1 4 g_2(L)$ und $h_3 = \frac 1 4 g_3(L)$.
Zunächst beweisen wir $m\leq 4$ als Induktionsanfang:}
\begin{enumerate}[leftmargin=*, itemsep=1ex]
    \item \en{From its definition in Lemma~\ref{\en{EN}lemdefifm}, we observe $F_1(z)=1$. Thus $F_1(z)=P_1(\wp(z))$.}%
    \de{In der Definition in Lemma~\ref{\en{EN}lemdefifm} lesen wir $F_1(z)=1$ ab. Also gilt $F_1(z)=P_1(\wp(z))$.}
    
    \item \en{In Lemma~\ref{\en{EN}lemdefifm} we proved that
    $F_2(z)=\frac{\sigma(2z)}{\sigma(z)^4}$ has a pole of order three at $z=0$ and three zeros at $z\in \DIV(2)$.
    Prop.~\ref{\en{EN}zerowp} shows that the same is true for $\wp'(z)$, thus $\frac{F_2(z)}{\wp'(z)}$ is constant by the first Liouville theorem (Prop.~\ref{\en{EN}liouville1}).
    Comparing the Laurent series at $z=0$ gives $F_2(z)=2z^{-3}+O(z^{-1})$ and $\wp'(z)=-2z^{-3} + O(z^{-1})$, thus $\frac{F_2(z)}{\wp'(z)}=-1$ and $F_2(z)=-\wp'(z)\cdot P_2(\wp(z))$.}%
    \de{In Lemma~\ref{\en{EN}lemdefifm} haben wir bewiesen, dass
    $F_2(z)=\frac{\sigma(2z)}{\sigma(z)^4}$ einen Pol dritter Ordnung bei $z=0$ and drei Nullstellen bei $z\in \DIV(2)$ hat.
    Satz~\ref{\en{EN}zerowp} besagt, dass dies auch auf $\wp'(z)$ zutrifft, also ist $\frac{F_2(z)}{\wp'(z)}$ aufgrund des ersten Liouville'schen Satzes~\ref{\en{EN}liouville1} konstant.
    Ein Vergleich der Laurentreihen um $z=0$ liefert $F_2(z)=2z^{-3}+O(z^{-1})$ und $\wp'(z)=-2z^{-3} + O(z^{-1})$, also $\frac{F_2(z)}{\wp'(z)}=-1$ und $F_2(z)=-\wp'(z)\cdot P_2(\wp(z))$.}
    
    \item \en{From Lemma~\ref{\en{EN}lemwpnz} with $n=2$ we obtain}%
    \de{Aus Lemma~\ref{\en{EN}lemwpnz} folgt mit $n=2$:}
    $\wp(2z)=\wp(z)-\frac{F_1(z)\cdot F_3(z)}{F_2(z)^2}$.
    \en{Here we use $F_1(z)=1$ and $F_2(z)=-\wp'(z)$ and deduce}%
    \de{Hier nutzen wir $F_1(z)=1$ und $F_2(z)=-\wp'(z)$ und erhalten:}
    $$F_3(z)=\wp'(z)^2\cdot\left(\wp(z)-\wp(2z)\right)$$
    \en{Here we use Lemma~\ref{\en{EN}lemwp2str} for $\wp(2z)$ and obtain}%
    \de{Hier setzen wir $\wp(2z)$ aus Lemma~\ref{\en{EN}lemwp2str} ein:}
    \begin{align*}
        F_3(z) &= \wp'(z)^2\cdot\left(\wp(z)-\frac 1 4 \left(\frac{\wp''(z)}{\wp'(z)}\right)^2+ 2 \wp(z)\right)
        = -\frac 1 4 \wp''(z)^2 + 3 \wp(z)\wp'(z)^2
    \end{align*}
    \en{Here we use the representation of $\wp''$ from Lemma~\ref{\en{EN}lemwp2str} and the algebraic differential equation from Prop.~\ref{\en{EN}dglP}:}%
    \de{Mit der Darstellung von $\wp''$ aus Lemma~\ref{\en{EN}lemwp2str} und der algebraischen Differentialgleichung der $\wp$-Funktion in Satz~\ref{\en{EN}dglP} folgt:}
    \begin{align*}
        F_3(z)
        &= -\frac 1 4 (6\wp(z)^2-2h_2)^2 + 3 \wp(z)\left(4\wp(z)^3-4h_2\wp(z)-4 h_3\right)
    \end{align*}
    \en{Using the abbreviation $x=\wp(z)$ yields}%
    \de{Mit der Abkürzung $x=\wp(z)$ führt das auf}
    \begin{align*}
        F_3(z)
        &= -\frac 1 4 (6x^2-2h_2)^2 + 3 x\left(4x^3-4h_2x-4 h_3\right)\\
        &= -9x^4+6h_2x^2-h_2^2+12x^4-12 h_2 x^2- 12 h_3 x = P_3(x)=P_3(\wp(z))
    \end{align*}
    
    \item \en{Again from the definition of the $F_m(z)$ in Lemma~\ref{\en{EN}lemdefifm} we obtain}%
    \de{Wir nutzen wieder die Definition der $F_m(z)$ aus Lemma~\ref{\en{EN}lemdefifm} und erhalten}
    $$F_4(z) = \frac{\sigma(4z)}{\sigma(z)^{16}}=\frac{\sigma(4z)}{\sigma(2z)^4}\cdot\left(\frac{\sigma(2z)}{\sigma(z)^4}\right)^4 = F_2(2z)\cdot F_2(z)^4$$
    \en{Using $F_2(z)=-\wp'(z)$ we deduce}%
    \de{Mit $F_2(z)=-\wp'(z)$ folgt}
    $$F_4(z) = -\wp'(2z)\cdot\wp'(z)^4$$
    \en{Next we write $\wp(2z)$ from Lemma~\ref{\en{EN}lemwp2str} only in terms of $x=\wp(z)$, $h_2 = \frac 1 4 g_2$ and $h_3 = \frac 1 4 g_3$:}%
    \de{Als nächstes drücken wir $\wp(2z)$ aus Lemma~\ref{\en{EN}lemwp2str} nur mit Hilfe von  $x=\wp(z)$, $h_2 = \frac 1 4 g_2$ und $h_3 = \frac 1 4 g_3$ aus:}
    \begin{align*}
    \wp(2z) &= \frac{1}{4}\cdot\frac{(6x^2-2h_2)^2}{4x^3-4h_2x-4h_3}-2x
    = \frac{9x^4-6h_2x^2+h_2^2}{4x^3-4h_2x-4h_3}-2x\\
    &= \frac{9x^4-6h_2x^2+h_2^2-2x(4x^3-4h_2x-4h_3)}{4x^3-4h_2x-4h_3}\\
    &= \frac{x^4+2h_2x^2+8h_3x+h_2^2}{4x^3-4h_2x-4h_3}=:f(x)
    \end{align*}
    \en{Deriving both sides with respect to $z$ yields $2\wp'(2z) = f'(x)\cdot\wp'(z)$, where}%
    \de{Beide Seiten nach $z$ ableiten liefert $2\wp'(2z) = f'(x)\cdot\wp'(z)$, wobei}
    \begin{align*}
        f'(x) &= \frac{\left(\begin{aligned}(4x^3+4h_2x+8h_3)\cdot(4x^3-4h_2x-4h_3)\\-(x^4+2h_2x^2+8h_3x+h_2^2)\cdot(12x^2-4h_2)\end{aligned}\right)}{(4x^3-4h_2x-4h_3)^2}\\
        &= \frac{4x^6-20h_2x^4-80h_3x^3-20h_2^2x^2-16h_2h_3x-32h_3^2+4h_2^3}{(4x^3-g_2x-g_3)^2}
        = \frac{2 P_4(x)}{\wp'(z)^4}
    \end{align*}
    \en{This yields}\de{Das liefert}
    $F_4(z) = -\wp'(2z)\cdot\wp'(z)^4 = - \frac 1 2 f'(x)\cdot\wp'(z)\cdot\wp'(z)^4 = -\wp'(z)\cdot P_4(x)$.
\end{enumerate}
\en{Now we have proven that the Proposition is true for $m\leq 4$.
To prove it for any $m\geq 5$, we assume that it is correct for all numbers less then $m$, and we distinguish four cases:}%
\de{Somit ist der Satz für $m\leq 4$ bewiesen.
Um ihn für $m\geq 5$ zu beweisen, nehmen wir die Korrektheit für alle Zahlen kleiner $m$ an und unterscheiden vier Fälle:}
\begin{enumerate}[leftmargin=*,itemsep=1ex]
    \item \en{If $m=4k+1$ with an integer $k\geq 1$, then Lemma~\ref{\en{EN}lemfrekur} tells (using $n=2k$):}%
    \de{Falls $m=4k+1$ ist mit natürlichem $k\geq 1$, liefert Lemma~\ref{\en{EN}lemfrekur} mit $n=2k$:}
            \begin{align*}
                F_{4k+1}(z) &= F_{2k+2}(z)\cdot F_{2k}(z)^3 - F_{2k-1}(z)\cdot F_{2k+1}(z)^3
            \end{align*}
            \en{Now, by induction, we get}\de{Aus der Induktionsvoraussetzung folgt:}
            \begin{align*}
                F_{4k+1}(z) &= -\wp'(z)P_{2k+2}(x)\cdot (-\wp'(z)P_{2k}(x))^3 - P_{2k-1}(x)\cdot P_{2k+1}(x)^3\\
                        &= \wp'(z)^4\cdot P_{2k+2}(x)\cdot P_{2k}(x)^3 - P_{2k-1}(x)\cdot P_{2k+1}(x)^3
            \end{align*}
            \en{Using}\de{Mit} $\wp'(z)^2=4\wp(z)^3-g_2\wp(z)-g_3 = 4(x^3-h_2x-h_3)$ \en{yields}\de{folgt}
            \begin{align*}
                F_{4k+1}(z) &= 16(x^3-h_2x-h_3)^2\cdot P_{2k+2}\cdot P_{2k}^3 - P_{2k-1}\cdot P_{2k+1}^3 = P_{4k+1}
            \end{align*}
            \en{which proves the Proposition in this case.}%
            \de{womit der Satz in diesem Fall bewiesen ist.}
    \item \en{If $m=4k+2$ with an integer $k\geq 1$, then Lemma~\ref{\en{EN}lemfrekur} tells (using $n=2k+1$):}%
    \de{Falls $m=4k+2$ mit natürlichem $k\geq 1$, liefert Lemma~\ref{\en{EN}lemfrekur} mit $n=2k+1$:}
            \begin{align*}
                F_{4k+2}\cdot F_2 &=  F_{2k+1}\cdot \left( F_{2k+3}\cdot F_{2k}^2 - F_{2k-1}\cdot F_{2k+2}^2\right)
            \end{align*}
            \en{Now, by induction, we get}\de{Aus der Induktionsvoraussetzung folgt:}
            \begin{align*}
                F_{4k+2}\cdot F_2 &=  P_{2k+1}\cdot \left( P_{2k+3}\cdot \wp'(z)^2 P_{2k}^2 - P_{2k-1}\cdot\wp'(z)^2 P_{2k+2}^2\right)
            \end{align*}
            \en{Dividing by $F_2=-\wp'(z)$ yields}\de{Hier dividieren wir durch $F_2=-\wp'(z)$ und erhalten:}
            \begin{align*}
                F_{4k+2} &=  -\wp'(z)\cdot P_{2k+1}\cdot \left( P_{2k+3}\cdot P_{2k}^2 - P_{2k-1}\cdot P_{2k+2}^2\right) = -\wp'(z)\cdot P_{4k+2}
            \end{align*}
            \en{Thus the Proposition is also true in this case.}\de{Somit ist der Satz auch in diesem Fall bewiesen.}
    \item \en{If $m=4k+3$ with an integer $k\geq 1$, then Lemma~\ref{\en{EN}lemfrekur} tells (using $n=2k+1$):}%
    \de{Falls $m=4k+3$ ist mit natürlichem $k\geq 1$, liefert Lemma~\ref{\en{EN}lemfrekur} mit $n=2k+1$:}
            \begin{align*}
                F_{4k+3} &= F_{2k+3}\cdot F_{2k+1}^3 - F_{2k}\cdot F_{2k+2}^3
            \end{align*}
           \en{Now, by induction, we get}\de{Aus der Induktionsvoraussetzung folgt:}
            \begin{align*}
                F_{4k+3} &= P_{2k+3}\cdot P_{2k+1}^3 - (-\wp'(z)P_{2k})\cdot(-\wp'(z) P_{2k+2})^3\\
                        &= P_{2k+3}\cdot P_{2k+1}^3 - \wp'(z)^4 P_{2k}\cdot P_{2k+2}^3
            \end{align*}
            \en{Again,}\de{Mit} $\wp'(z)^2=4\wp(z)^3-g_2\wp(z)-g_3 = 4(x^3-h_2x-h_3)$ \en{yields}\de{folgt wiederum}
             \begin{align*}
                F_{4k+3} = P_{2k+3}\cdot P_{2k+1}^3 - 16(x^3-h_2x-h_3)^2 P_{2k}\cdot P_{2k+2}^3 = P_{4k+3}
            \end{align*}
             \en{which proves the Proposition in this case.}%
            \de{womit der Satz auch in diesem Fall bewiesen ist.}
    \item \en{If $m=4k+4$ with an integer $k\geq 1$, then Lemma~\ref{\en{EN}lemfrekur} tells (using $n=2k+2$):}%
    \de{Falls $m=4k+4$ mit natürlichem $k\geq 1$, liefert Lemma~\ref{\en{EN}lemfrekur} mit $n=2k+2$:}
            \begin{align*}
                F_{4k+4}\cdot F_2 &=  F_{2k+2}\cdot \left( F_{2k+4}\cdot F_{2k+1}^2 - F_{2k}\cdot F_{2k+3}^2\right)
            \end{align*}
            \en{Now, by induction, we get}\de{Aus der Induktionsvoraussetzung folgt:}
            \begin{align*}
                F_{4k+4}\cdot F_2 &=  -\wp'(z) P_{2k+2}\cdot \left( -\wp'(z) P_{2k+4}\cdot P_{2k+1}^2 - (-\wp'(z)P_{2k})\cdot P_{2k+3}^2\right)
            \end{align*}
            \en{Dividing by $F_2=-\wp'(z)$ yields}\de{Hier dividieren wir durch $F_2=-\wp'(z)$ und erhalten:}
            \begin{align*}
                F_{4k+4} &=  -\wp'(z)\cdot P_{2k+2}\cdot \left( P_{2k+4}\cdot P_{2k+1}^2 - P_{2k}\cdot P_{2k+3}^2\right) = -\wp'(z)\cdot P_{4k+4}
            \end{align*}
            \en{Thus the Proposition is also true in this case.}\de{Somit ist der Satz auch in diesem Fall bewiesen.}
\end{enumerate}
\en{By induction, we have proven Prop.~\ref{\en{EN}propfmpm}.}%
\de{Satz~\ref{\en{EN}propfmpm} ist also für alle natürlichen $m\neq 0$ per vollständiger Induktion bewiesen.}
\end{proof}


\begin{theo}\label{\en{EN}thmbaker}
\en{For any positive integer $m$, it holds:}%
\de{Für alle natürlichen $m\neq 0$ gilt}
$$m^2\cdot\prod_{u\in \DIV(m)}\left(x-\wp(u)\right) = \cas{4\cdot (x^3-h_2x - h_3) \cdot P_m^2(x)}{P_m^2(x)}$$
\end{theo}
\begin{proof}
\en{We will prove that this is true for every $x\in \mathbb C$:}%
\de{Wir werden diese Aussage für alle $x\in \mathbb C$ beweisen:}
\en{Given $x\in\mathbb C$. Then chose $z\in\mathbb C$ such that $\wp(z)=x$. This is possible, since the $\wp$-function takes on every value (apply the third Liouville theorem~\ref{\en{EN}liouville3} to $f(u):=\wp(u)-x$ for the given $x$ and chose $z$ as one of the zeros of $f(u)$).}%
\de{Sei also ein beliebiges $x\in \mathbb C$ gegeben. Wählen dann ein $z\in\mathbb C$ mit $\wp(z)=x$. Das ist möglich, weil die $\wp$-Funktion jeden Wert annimmt (man wende den dritten Liouville'schen Satz~\ref{\en{EN}liouville3} auf $f(u):=\wp(u)-x$ an und wähle für $z$ eine der Nullstellen von $f(u)$).}
\en{Then Prop.~\ref{\en{EN}propfmprod} tells us that the left hand side is equal to $F_m(z)^2$.
Prop.~\ref{\en{EN}propfmpm} gives}%
\de{Aus Satz~\ref{\en{EN}propfmprod} folgt, dass die linke Seite gleich $F_m(z)^2$ ist.
Satz~\ref{\en{EN}propfmpm} liefert dann}
$$F_m(z)^2 = \cas{(-\wp'(z))^2 \cdot P_m^2(\wp(z))}{P_m^2(\wp(z))}$$
\en{Prop.~\ref{\en{EN}dglP} tells $(-\wp'(z))^2=4\wp(z)^3-g_2\wp(z)-g_3=4(x^3-h_2x-h_3)$, thus we have proven Thm.~\ref{\en{EN}thmbaker} for the given $x\in\mathbb C$.
This proves the Theorem, because $x$ was chosen arbitrarily.}%
\de{Satz~\ref{\en{EN}dglP} besagt $(-\wp'(z))^2=4\wp(z)^3-g_2\wp(z)-g_3=4(x^3-h_2x-h_3)$, also ist Thm.~\ref{\en{EN}thmbaker} für das gegebene $x\in\mathbb C$ und somit für alle $x$ bewiesen.}
\end{proof}

\begin{thm}\label{\en{EN}propindu}
\en{For any positive integer $m$, it holds:}%
\de{Für alle natürlichen $m\neq 0$ gilt}
\begin{enumerate}
    \item \en{$P_m$ is a polynomial in $x$, $h_2$ and $h_3$ with coefficients in $\mathbb Z$.}%
    \de{$P_m$ ist ein Polynom in $x$, $h_2$ und $h_3$ mit Koeffizienten aus $\mathbb Z$.}\\[2ex]
    \hspace*{-1.15cm}\en{Furthermore, when $P_m$ is regarded as a polynomial in $x$, it holds:}%
    \de{Weiter gilt, wenn man $P_m$ als Polynom in $x$ betrachtet:}
    \item \en{The degree of $P_m(x)$ is}\de{Der Grad von $P_m(x)$ ist} $d_m:=\cas{\frac{m^2-4}{2}}{\frac{m^2-1}{2}}$.
    \item \en{The leading coefficient of $P_m(x)$ is $\pm l_m$ with}\de{Der Leitkoeffizient von $P_m(x)$ ist $\pm l_m$ mit} $l_m:=\cas{\frac{m}{2}}{m}$.
    \item \en{The second highest coefficient of $P_m(x)$ is $0$.}%
    \de{Der zweithöchste Koeffizient von $P_m(x)$ ist $0$.}
\end{enumerate}
\end{thm}
\begin{proof}
\begin{enumerate}[leftmargin=*,itemsep=0.5ex]
    \item \en{Since the Definition~\ref{\en{EN}defipm} of $P_m$ only includes multiplication and summation, statement (1) is correct for all $m$ by induction over $m$.}%
        \de{Da die Definition~\ref{\en{EN}defipm} der $P_m$ nur Multiplikationen und Additionen enthält, folgt Aussage (1) durch vollständige Induktion über $m$.}
    \item \en{Lemma~\ref{\en{EN}lemmm} tells that there are $m^2-1$ values $u\in \DIV(m)$, thus the left-hand-side of Thm.~\ref{\en{EN}thmbaker} has the degree $m^2-1$. This proves statement (2).}%
        \de{Lemma~\ref{\en{EN}lemmm} besagt, dass es $m^2-1$ Werte $u\in \DIV(m)$ gibt, also hat die linke Seite in Thm.~\ref{\en{EN}thmbaker} den Grad $m^2-1$. Hieraus folgt Aussage (2).}
    \item \en{On the left-hand-side of Thm.~\ref{\en{EN}thmbaker}, we observe that the leading coefficient is $m^2$, thus the leading coefficient of $P_m(x)$ has to be either $l_m$ or $-l_m$.
    It could be proven by induction over $m$ that it is indeed $l_m$, but for our means it's enough to prove statement (3) up to a factor of $\pm 1$.}%
    \de{Auf der linken Seite von Thm.~\ref{\en{EN}thmbaker} erkennen wir den Leitkoeffizient $m^2$, also muss der Leitkoeffizient von $P_m(x)$ entweder $l_m$ oder $-l_m$ sein.
    Man könnte per vollständiger Induktion beweisen, dass es tatsächlich $l_m$ ist, aber für unsere Zwecke ist es ausreichend, Aussage (3) bis auf einen Faktor von $\pm 1$ zu beweisen.}
    \item \en{For proving (4) we use the fact that for any two polynomials $f(x)=\sum_{k=0}^{n} a_k x^k$ and $g(x)=\sum_{k=0}^{m} b_k x^k$ it holds:
The second highest coefficient of $f(x)\cdot g(x)$ is given by $a_n\cdot b_{m-1}+a_{n-1}\cdot b_m$.
This proves that if the second highest coefficients of the factors vanish, the second highest coefficient of the product also vanishes.
Further we observe that in the recursive definition of the $P_m$, we only add or substract polynomials of the same degree (cf.~(3)), which proves (4) by induction.}%
\de{Um (4) zu beweisen, nutzen wir die Tatsache, dass für zwei beliebige Polynome $f(x)=\sum_{k=0}^{n} a_k x^k$ und $g(x)=\sum_{k=0}^{m} b_k x^k$ gilt:
Der zweithöchste Koeffizient von $f(x)\cdot g(x)$ ist $a_n\cdot b_{m-1}+a_{n-1}\cdot b_m$.
Also: Wenn die zweithöchsten Koeffizienten der Faktoren Null sind, dann ist auch der zweithöchste Koeffizient des Produktes Null.
Des weiteren bemerken wir, dass in der rekursiven Definition der $P_m$ nur Polynome gleichen Grades addiert oder subtrahiert werden (vgl.~(3)), woraus (4) per Induktion folgt.}
\end{enumerate}\vspace*{-12pt}
\end{proof}

\begin{theo}\label{\en{EN}theowpsum}
\en{For any positive integer $m$, the following sum of $m$-division-values vanishes:}%
\de{Für alle natürlichen $m\neq 0$ verschwindet die folgende Summe von $m$-Teilungswerten:}
$$\sum_{u\in \DIV(m)}\wp(u)=0$$
\end{theo}
\begin{proof}
\en{Theorem~\ref{\en{EN}thmbaker} gives two equivalent representations of a polynomial in $x$.
We multiply out the left hand side of Thm.~\ref{\en{EN}thmbaker} and obtain}%
\de{Theorem~\ref{\en{EN}thmbaker} liefert zwei äquivalente Darstellungen eines Polynoms in $x$.
Wir multiplizieren die linke Seite des Thm.~\ref{\en{EN}thmbaker} aus und erhalten}
$$m^2\cdot x^{m^2-1} - \left(m^2\cdot\sum_{u\in \DIV(m)} \wp(u)\right)\cdot x^{m^2-2} + \sum_{k=0}^{m^2-3} a_k x^k$$
\en{Now Prop.~\ref{\en{EN}propindu} (4) tells us that the second highest coefficient of the right hand side of Thm.~\ref{\en{EN}thmbaker} is zero -- thus the same coefficient on the left hand side also vanishes. This proves Thm.~\ref{\en{EN}theowpsum}.}%
\de{Jetzt besagt Satz~\ref{\en{EN}propindu} (4), dass der zweithöchste Koeffizient auf der rechten Seite von Thm.~\ref{\en{EN}thmbaker} Null ist -- also muss dieser Koeffizient auch auf der linken Seite verschwinden. Dies beweist Thm.~\ref{\en{EN}theowpsum}.}
\end{proof}


\begin{theo}\label{\en{EN}theomwpu}
\en{Given a lattice $L$, let $h_2:=\frac{1}{4}g_2(L)$ and $h_3:=\frac{1}{4}g_3(L)$.
Then for any positive integer $m$ and for any $m$-division-point $u\in \DIV(m)$, the term}%
\de{Sei $L$ ein Gitter, sei $h_2:=\frac{1}{4}g_2(L)$ und $h_3:=\frac{1}{4}g_3(L)$.
Dann gilt für alle natürlichen $m\neq 0$ und für alle $m$-Teilungsstellen $u\in \DIV(m)$, dass}
$$m\cdot\wp(u)$$
\en{is an algebraic integer of~~$\mathbb Z[h_2;h_3]$.}%
\de{ganzalgebraisch in $\mathbb Z[h_2;h_3]$ ist.}\\
\en{Moreover, if $m$ is even, then $\frac{m}{2}\cdot\wp(u)$ is also an algebraic integer of~~$\mathbb Z[h_2;h_3]$.}%
\de{Und falls $m$ gerade ist, ist $\frac{m}{2}\cdot\wp(u)$ ebenfalls ganzalgebraisch in $\mathbb Z[h_2;h_3]$.}
\end{theo}
\begin{proof}
\en{Throughout this proof, we denote $\mathbb I := \mathbb Z[h_2;h_3]$.
From Thm.~\ref{\en{EN}thmbaker} we know that the $\wp(u)$ with $u\in \DIV(m)$ are either among the zeros of $x^3-h_2x-h_3$ or among the zeros of $P_m(x)$.

The zeros of $x^3-h_2x-h_3$ are algebraic integers of $\mathbb I$, so these $\wp(u)$ are algebraic integers of $\mathbb I$.
Since $P_1=P_2=1$, it remains to check the zeros of $P_m(x)$ for $m\geq 3$.}%
\de{Für diesen Beweis bezeichnen wir $\mathbb I := \mathbb Z[h_2;h_3]$.
Aus Thm.~\ref{\en{EN}thmbaker} wissen wir, dass die $\wp(u)$ mit $u\in \DIV(m)$ entweder Nullstellen von $x^3-h_2x-h_3$ oder von $P_m(x)$ sind.

Die Nullstellen von $x^3-h_2x-h_3$ sind ganzalgebraisch in $\mathbb I$, also sind diese $\wp(u)$ ganzalgebraisch in $\mathbb I$.
Wegen $P_1=P_2=1$ müssen wir nur noch die Nullstellen von $P_m(x)$ für $m\geq 3$ prüfen.}

\en{Again, we use the notation $d_m:=\cas{\frac{m^2-4}{2}}{\frac{m^2-1}{2}}$ for the degree of $P_m$ (as in Prop.~\ref{\en{EN}propindu}) and the notation $l_m:=\cas{\frac{m}{2}}{m}$.}%
\de{Hier nutzen wir wieder die Notation $d_m:=\cas{\frac{m^2-4}{2}}{\frac{m^2-1}{2}}$ für den Grad von $P_m$ (vgl.~Satz~\ref{\en{EN}propindu}) und die Notation $l_m:=\cas{\frac{m}{2}}{m}$.}

\en{Then we define the polynomial $\displaystyle h(x):=P_m\lk\frac{x}{l_m}\rk\cdot l_m^{d_m-1}$.
From Prop.~\ref{\en{EN}propindu} (3) we know that the leading coefficient of $P_m(x)$ is either $l_m$ or $-l_m$. This leads to}%
\de{Dann betrachten wir das Polynom $\displaystyle h(x):=P_m\lk\frac{x}{l_m}\rk\cdot l_m^{d_m-1}$.
Da der Leitkoeffizient von $P_m(x)$ nach Satz~\ref{\en{EN}propindu} (3) entweder $l_m$ oder $-l_m$ ist, erhalten wir}
\begin{align*}
h(x) &= \pm l_m \cdot \left(\frac{x}{l_m}\right)^{d_m}\cdot \left(l_m\right)^{d_m-1} + \sum_{k=0}^{d_m-2}b_k \left(\frac{x}{l_m}\right)^{k}\cdot \left(l_m\right)^{d_m-1}\\
&= \pm x^{d_m} + \sum_{k=0}^{d_m-2}b_k \cdot\left(l_m\right)^{d_m-1-k}\cdot x^k
\end{align*}
\en{where the $b_k$ are the coefficients of $P_m(x)$, for which Prop.~\ref{\en{EN}propindu}~(1) tells: $b_k\in\mathbb I$.}%
\de{wobei $b_k$ die Koeffizienten von $P_m(x)$ sind, für die nach Satz~\ref{\en{EN}propindu} (1) gilt: $b_k\in\mathbb I$.}

\en{This shows that $h$ is a monic polynomial in $x$ whose coefficients are in $\mathbb I$, so the zeros of $h(x)$ are algebraic integers of $\mathbb I$.
But it holds: $$h(x)=0\quad\Longleftrightarrow\quad P_m\lk\frac x {l_m}\rk=0$$
This proves that for all $\wp(u)$ which are zeros of $P_m(x)$, the term $l_m\cdot \wp(u)$ is an algebraic integer of $\mathbb I$.
Altogether, we have proved Thm.~\ref{\en{EN}theomwpu} for all $u\in \DIV(m)$.}%
\de{Dies zeigt, dass $h$ ein monisches Polynom in $x$ ist, dessen Koeffizienten in $\mathbb I$ sind -- also sind die Nullstellen von $h(x)$ ganzalgebraisch in $\mathbb I$.
Aber es gilt: $$h(x)=0\quad\Longleftrightarrow\quad P_m\lk\frac x {l_m}\rk=0$$
Hieraus folgt auch für alle $\wp(u)$, die Nullstellen von $P_m(x)$ sind, dass $l_m\cdot \wp(u)$ ganzalgebraisch in $\mathbb I$ ist.
Insgesamt haben wir Thm.~\ref{\en{EN}theomwpu} also für alle $u\in \DIV(m)$ bewiesen.}
\end{proof}